\documentclass[CJK]{paper}  %aspectratio控制页面的宽高比，16:9充满屏幕
%\documentclass[CJK,aspectratio=43]{beamer} %4:3宽高比

\usepackage{ctex} %中文字体设置
\usepackage{hyperref}
\usepackage{multicol}
\usepackage{makecell}

\usepackage{graphicx}
\usepackage{float}
\usepackage{subfigure}

\usepackage{amsmath}
\usepackage{mathtools}

\usepackage{soul}

\usepackage{diagbox}

\usepackage{tcolorbox}


% \usepackage[style=alphabetic,maxnames=3,minnames=1,maxbibnames=99]{biblatex}
% \addbibresource{ref.bib} 


\begin{document}



\title{Details on Parameterized Tests of Gravity Theories Using GWs with Matched Filtering}
% \subtitle{子题目}
% \ctexset{today=old}
% \author{ldc}
\date{today}
% 生成标题
\maketitle


\newpage
\tableofcontents


\newpage
\section{检验类型}

\subsection{相容性检验}
\par 不同gw检验间的关系：\href{https://arxiv.org/abs/2109.06988}{Investigating the relation between gravitational wave tests of general relativity} \cite{gw_tests_relation}


\subsubsection{Residual Test}

Residual Test的优点：
\begin{itemize}
	\item[1. ] 不与双星参数简并\cite{test_GR_GW150914}；
\end{itemize}


\subsubsection{IMR Test}


\subsubsection{QNM Test}


\subsubsection{Parameterized Tests}

\par Deviations 类型：
\begin{table}[h!]
	\begin{center}
		\begin{tabular}{|c|c|}
			\hline
			\textsc{Type} & 表达式 \\
			\hline
			Fractional & $p_i \longrightarrow p_i \cdot (1 + \delta \hat{p}_i)$ \\
			\hline
			Absolute/Additional & $p_i \longrightarrow p_i + \delta p_i$ \\
			\hline
			Particular Functional Form &  $p_i \longrightarrow p_i \cdot (1 + \delta p_i(\Xi;f))$\\
			\hline
		\end{tabular}
	\end{center}
\end{table}


\subsubsection{Dispersion Relation Test}


\subsubsection{Polarization Test}


\subsubsection{$\chi^2$ GW Search}

\subsection{测量特定理论参数}

\par 测量特定理论的参数，原则上必须采用该理论全阶段non-GR波形，然而困境是\cite{wf_modelling_for_LISA}：
\begin{itemize}
	\item [1. ] 理论结果；
	\item [2. ] if a well-posed initial value problem；Numerical Results；
\end{itemize}

\par 折中方案：
\begin{itemize}
	\item[1. ] 截取CBC波形模板中已有明确理论解的特定阶段；
	\item[2. ] 可能的方案是 ppE\cite{ppE}；
	\item[3. ] 不管，全阶段、但是只改一个阶段\cite{test_GR_GWTC3}；
	尽管\cite{test_GR_GWTC3}前的每篇test GR都提到如果要constrain某个特定理论参数、必须要用完整的NON-GR模板，但这里：
	\begin{quote}
		Likewise, one would also need a robust estimate of 
		the error caused by neglecting currently unknown higher-order
		PN corrections and \emph{deviations in the merger–ringdown phase},
		which will also differ from the GR one [195, 196].
	\end{quote}
	\par 此外考虑的点有，这里提到的结果对参数的上界估计依赖具体的参数化(如NON-GR修正的截断频率、单/多参数)、先验分布、适定性(well-possedness)。
\end{itemize}



\newpage
\section{GR波形模板误差}

\subsection{ppE-type deviations innaccuracy due to HM和Precessing}
\par 实际上详细论述了参数估计或参数化检验中可能遇到的各种包括模板在内的系统误差：\cite{systematic_bias_precession_HM}

\subsection{具体考虑}
\begin{itemize}
	\item[1. ] GW150914\cite{test_GR_GW150914}: 
		
		\par SEOBNRv2\_ROM\_DOUBLESPIN: double-spin, non-precessing；
		\par IMRPhenomPv2: single-effective-spin, precessing；
		\par NR validation for Residual Test: we added the publicly available and new NR waveforms as mock signals to the data in the neighbourhood of GW150914 [36, 50, 51]；
		\par NR validation for Parameterized Test: 用一个 GW150914-like 的NR波形
			注入到不同 noise-segments 的20个数据集，其中一个案例复现接近、其它情况偏离比信号情况小得多尽管宽度comparable；
			因此单参分析中，不认为 instrumental noise fluctuations会导致自由度的明显偏离even in the absence of an actual deviation from GR，但不排除系统误差或者模板中有physics未被考虑；
		\par Detector Calibration: These uncertainties are modeled and included in the results presented here according to the treatment detailed in Ref. [3].
	\item[2. ] GW170817\cite{test_GR_GW170817}：
		\par NRTidal: PN inspiral including finite size effects；alse merger and ringdown；tidal effects。
		\par 往BBH波形 phasing 上加一个 $\phi_\text{T}(f)$(通过混合PN信息、tidal EOB模型以及NR input来build) 变成 Tidal 波形；IMRPhenomPNRT；
		\par windowing the amplitude such that the merger and ringdown are smoothly removed from the model;
			Because of the unknown merger-ringdown behavior in the case of binary neutron stars, which in any case gets removed from the waveform model,
		\par 对 IMRPhenomPv3 进行以上流程，得到 PhenomPNRT；
		\par SEOBNRT, augmented from SEOBNRv4；
		\par 两个波形的连续性条件选择不同、但是由于post-inspiral截断，影响认为不大；quadrupole-monopole couplings的一些计算选择；
	\item[] 
\end{itemize}



\newpage
\section{波形模板构造}

\subsection{参数化策略}
\subsubsection{Tiger}
\par \cite{TIGER_towards}\cite{TIGER_towards_further}\cite{TIGER}

\subsubsection{ppE}
\begin{itemize}
	\item[1. ] ppE for equal-mass, non-spin BBH: \cite{ppE}；
	\item[2. ] ppE for precessing binaries: \cite{ppE_precessing}；
\end{itemize}

\subsubsection{FTI}
\par \cite{FTI}

\subsubsection{pSEOBNR}
\begin{itemize}
	\item[1. ] pSEOBNRv4HM\cite{pSEOBNRv4HM}；
	\item[2. ] pSEOBNRv5HM\cite{pSEOBNRv5HM}；
	\item[3. ] parameterized Plunge-Merger-Ringdown\cite{pPMR}；
\end{itemize}

\subsubsection{pyRing}


\subsubsection{ParSpec}
\begin{itemize}
	\item[1. ] ParSpec\cite{ParSpec}；
	\item[2. ] High Spin ParSpec\cite{ParSpec_high_spin}；
\end{itemize}


\subsection{LALSimulation的IMRPhenom默认参数化}

\subsubsection{Parameterized PN Inspiral}

\begin{tcolorbox}
	\begin{align}
		\Phi(f) = 2\pi ft_c - \varphi_c - \frac{\pi}{4} + \Psi(f_k)
	\end{align}
	\begin{align}
		\Psi(f) &= \text{CCE}(k;\eta,M;f) \cdot \sum_j \left[ \psi_j + \psi_{jl} \ln{f} \right] f^{j/3}\\
				&= \text{CCE}(k;\eta,M;f) \cdot \sum_j \left[ \alpha_j + \alpha_{jl} \ln{(2\pi Mf)} \right] (2\pi Mf)^{j/3}
	\end{align}
	\begin{align}
		\text{CCE}(\eta,M;f) = k \cdot \frac{3}{256}\frac{1}{\eta}(2\pi Mf)^{-5/3}
	\end{align}
	\par 其中 $f_k = \frac{f}{k}$；大多数情况下只考虑 $k = 2$ 的四极矩情况.
\end{tcolorbox}

\begin{itemize}
	\item[1. ] PN系数的定义选择：

	\begin{table}[h!]
		\begin{center}
			\begin{tabular}{|c|c|c|c|}
				\hline
				\diagbox{} & 幂项 & 公共系数 & PN系数 \\
				\hline
				\cite{reexpress_pn_3} & $(\pi Mf)^{j/3}$ & & $\psi_j,\psi_{jl}$ \\
				\hline
				IMRPhenomD & $(\pi Mf)^{j/3}$ & & $\alpha_j,\alpha_{jl}$\\
				\hline
				IMRPhenomXAS & $(\pi Mf)^{j/3}$ & & $\alpha_j,\alpha_{jl}$\\
				\hline
				LALSimulation & $(Mf)^{j/3}$ & & \makecell{幂项的$\pi^{1/3}$乘进系数，\\ 调整$\alpha_j,\alpha_{jl}$} \\
				\hline
				ppE\cite{ppE} & $(\pi\mathcal{M}_c f)^{j/3}$ & $\text{CCE}(\eta,\mathcal{M}_c;f)$ & \makecell{$M$变成$\mathcal{M}_c$，\\ 调整$\alpha_j,\alpha_{jl}$} \\
				\hline
				GW150914 \cite{test_GR_GW150914} & \multicolumn{3}{|c|}{IMRPhenomD} \\
				\hline
				GWTC-3 \cite{test_GR_GWTC3} & $\pi \tilde{f} = \pi \frac{GM}{c^3}(1+z)f$ & & \\
				\hline
			\end{tabular}
		\end{center}
	\end{table}
	\item[2. ] IMRPhenom波形family有显式的PN展开形式，
			而没有者、比如SEOBNR family，对应的是一种更 generic 的PN系数 test 方案：
	\begin{table}[h]
		\begin{center}
			\begin{tabular}{|c|c|}
				\hline
				\textsc{Context} & PN系数的修改方法 \\
				\hline
				IMRPhenom & $\delta\varphi_j$  , Fractional or Additional(如$\delta\varphi_{-2}$，$\delta\varphi_1$) \\
				\hline
				SEOBNR & $\varphi_j \cdot \delta\varphi_j f^{(j-5)/3}$,log项同理\cite{test_GR_GW170817}\\
				\hline
			\end{tabular}
		\end{center}
	\end{table}
	\item[3. ] 是否将PN系数自旋依赖包括进偏离？
		\par 也就是考虑 $\varphi_j \longrightarrow (1 + \delta\varphi_j)\varphi_j^\text{NS} + \varphi_j^\text{S}$ \cite{test_GR_BBH_O1}；
		\par GW150914\cite{test_GR_GW150914}中包括了自旋依赖：At each point in parameter space,
			the coeffcients $p_i$ are evaluated for the local physical parameters (masses, spins)
			and multiplied by factors (1 + $\delta \hat{p}_i$).；
		\par GW170817\cite{test_GR_GW170817}中，\textbf{PhenomPNRT}的修正仅限于 non-spinning 部分
			(这个脚注2的说法和150914的说法矛盾了：There is no fundamental reason for this choice;
			we follow the convention used in previous publications on parameterized tests of GR [13, 15, 16, 27, 64].)
			(脚注3和；有一些讨论，详见论文)；
			\textbf{SEOBNRT}则是包括了自旋依赖；
		\par GWTC-3\cite{test_GR_GWTC3}中提到这种参数化可以避免系数偏离的 potentially singular behavior (奇异性原因来自non-spin和spin-dependent部分的 cancellation)。
	\item[4. ] 是否截断波形？
		\par 动力学情况不清楚的时候就要截断，比如用BNS事件GW170817 test GR\cite{test_GR_GW170817}。
		\par 或者利用某一特定阶段来约束特定 gravity theory。
	\item[5. ] 是否考虑 Amplitude 的参数化？
\end{itemize}


\subsubsection{Parameterized Intermediate}
IMRPhenomXAS：
\begin{align}
	\eta \varphi_\text{Int}^\prime = b_4f^{-4} + b_3f^{-3} + b_2f^{-2} + b_1f^{-1} + b_0 - \frac{4 c_0 a_\varphi}{(f-f_\text{RD})^2 + (2f_\text{damp})^2}
\end{align}

\subsubsection{Parameterized Ringdown}

IMRPhenomXAS (XAS中用 $p_i = -4, -2, -1/3$，D多一项$2$)：
\begin{align}
	\eta \varphi_\text{RD}^\prime = c_\text{RD} + \sum_i c_i f^{-p_i} + \frac{c_0 a_\varphi}{f_\text{damp}^2 + (f-f_\text{RD})^2}
\end{align}


\subsection{截断策略}
\begin{itemize}
	\item[1. ] GW150914
	\item[2. ] GW170817 \cite{test_GR_GW170817}
		\par windowing the amplitude such that the merger and ringdown are smoothly removed from the model;
\end{itemize}


\subsection{Single/Multiple at a time}
\begin{itemize}
	\item[1. ]\par \cite{test_GR_GW150914}：
		\par 考虑multiple是因为在大多数主流替代理论中，
		即便只考虑 leading order，PN inspiral相位演化系数每一PN阶都可能有理论修正；
		single的情况是期望capture predominant效应 (但是似乎并不然\cite{false_parametrization})；
		\par multiple的情况由于各阶符号导致的抵消，估计结果会 much more conservative；
		\par SVD\cite{PN_coeffcients_SVD}；
\end{itemize}




\newpage
\section{参数选择考量}

\par 不正确的参数化造成的影响：\cite{false_parametrization}

\subsection{SNR}

\subsubsection{Detector Sensitivity Band}
\subsubsection{事件IMR性质}
\begin{itemize}
	\item[1. ] GW150914\cite{test_GR_GW150914}
		\par For the intermediate and merger–ringdown regime, we
		verified that these intrinsic uncertainties are much smaller than
		the corresponding statistical uncertainties for GW150914, and
		thus do not affect our conclusions.
		\par  In the late-inspiral case,
		the uncertainties associated with the calibration of the j pa
		rameters are large, and almost comparable with the statistical
		measurement uncertainties. For this reason, we do not report
		results for the j parameters.
	\item[] 
\end{itemize}

\subsection{物理现象映射}
\begin{itemize}
	\item[1. ] GW150914\cite{test_GR_GW150914}
	\item[2. ] GW170817\cite{test_GR_GW170817}
		\par -1 PN的约束$\delta_{-2}$对应 dipole radiation。
	\item[] 
\end{itemize}


\newpage
\section{估计精确度}
\subsection{先验分布设置}
\begin{itemize}
	\item[1. ] GW150914\cite{test_GR_GW150914}
	\begin{table}[h!]
		\begin{center}
			\begin{tabular}{|c|c|}
				\hline
				\textsc{Deviation} & Prior Settings \\
				\hline
				$\delta \hat{\varphi}_i$ & Uniform$([-20,20])$ \\
				\hline
				$\delta \hat{\beta}_i$ & Uniform$([-3,3])$ \\
				\hline
				$\delta \hat{\alpha}_i$ & Uniform$([-5,5])$ \\
				\hline
			\end{tabular}
		\end{center}
	\end{table}
	\item[2. ] GW170817\cite{test_GR_GW170817}
		\par 由于$\delta\hat{\varphi}_0$与$\mathcal{M}_c$之间的 degeneracy，会比GR模板情况放宽$\mathcal{M}_c$；
		\par The conservative spin settings that $|\mathbf{\chi}_i| \leq 0.99$,
			allow the constraints on $\delta\hat{\varphi}_n$ to be directly compared with those from binary black hole observations,
			which used the same prior [13, 15]；
		\par $\Lambda_i < 5000$，多次后验分布tail触边；如果不强加这个约束、$\Lambda_i$与$\delta\hat{\varphi}_n$的correlation会导致后者的 weaker upper bounds；
\end{itemize}

\subsection{SNR}

\subsubsection{多事件}
\par \cite{combined_pi}\cite{combined_theory}\cite{combined_hierarchical}
\par Hiearchical Model with Hyper-parameters over Hyper-distributions 

\subsubsection{Detector Sensitivity Band}
\begin{itemize}
	\item[1. ] GW150914\cite{test_GR_GW150914}
		\par On the other hand, it is to be noted
		 that the LIGO detectors are still a factor of a few away from
		 their final design sensitivities [19]；
		\par \cite{LIGO_enhanced_sensitivity}\cite{AdvLIGO}，
		\par  As an aside, we note that GW150914 was detected with the
			LIGO detectors at about one-third of their final design sen
			sitivity, which is expected to be achieved around 2019 [19].
	\item[] 
\end{itemize}
\begin{table}[h!]
	\begin{center}
		\begin{tabular}{|c|c|}
			\hline
			事件 & 频带 span \\
			\hline
			GW150914 & \makecell{$\dot{P}_\text{orb} \sim -0.1$ at $f_\text{GW} \sim 30 \text{Hz}$\\ $\dot{P}_\text{orb} \sim -1$ at $f_\text{GW} \sim 132 \text{Hz}$\\ $20--300$ Hz\cite{test_GR_GW150914}}\\
			\hline
			GW170817 & \\
			\hline
		\end{tabular}
	\end{center}
\end{table}

\subsubsection{事件IMR性质}


\newpage
\section{事件性质}

\begin{tcolorbox}
	\par \href{https://waveview.cardiffgravity.org/}{Cardiff: Gravitational Waveform Viewer}；
	\par \href{https://peviewer.igwn.org/}{IGWN: GW Event Viewer}；
	\par \href{https://gwosc.org/eventapi/html/GWTC/?page=1&pagesize=100}{GWOSC: GWTC}；
\end{tcolorbox}

	\begin{table}[h!]
		\begin{center}
			\begin{tabular}{|c|c|}
				\hline
				事件 & 频带 span \\
				\hline
				GW150914 & \makecell{$\dot{P}_\text{orb} \sim -0.1$ at $f_\text{GW} \sim 30 \text{Hz}$\\ $\dot{P}_\text{orb} \sim -1$ at $f_\text{GW} \sim 132 \text{Hz}$\\ $20--300$ Hz\cite{test_GR_GW150914}}\\
				\hline
				GW170817 & \\
				\hline
			\end{tabular}
		\end{center}
	\end{table}


\newpage
\section{不同模板的后验样本混合策略}



\newpage
\section{GWs的约束结果}
\subsection{General Speakings}
\par \cite{test_GR_GWTC3}Likewise, one would also need a robust estimate of
 the error caused by neglecting currently unknown higher-order
 PNcorrections and deviations in the merger–ringdown phase,
 which will also di er from the GR one [195, 196]. 

\begin{itemize}
	\item[$\bullet$] GW150914\cite{test_GR_GW150914}
		\par 强引力场动力学的约束不强，如一些PN系数的相对偏离 $\sim \mathcal{O}(1)$
		\par LIGO detectors are still a factor of a few away from their final design sensitivities.
\end{itemize}


\subsubsection{Residual Test}


\subsubsection{IMR Test}

\begin{itemize}
	\item[1. ] GW150914\cite{test_GR_GW150914}: 
		\par $f_\text{GW}^\text{end insp} = 132 \text{Hz}$;
		\par Expected $\text{SNR}_\text{det} \sim 19.5$, $\sim 16$ if truncated between $132$ and $1024$ Hz.
		\par We have verified that these conclusions are not affected by the specific formula [40, 60, 63] used to predict $M_text{f}$ and $a_\text{f}$, nor affected by the choice of $f_\text{end}^\text{insp}$ GW within $50$ Hz.
	\item[2. ]
\end{itemize}


\subsubsection{QNM Test}


\subsubsection{Parameterized Tests}
\begin{itemize}
	\item[1. ] GW150914\cite{test_GR_GW150914}: 
		\par 对比J0737-3039， significantly informative except at 0PN order, where the double-pulsar
		boundis better thanks to the long observation time ( 10 years
		against $\sim 0.4$ s for GW150914). Thus, GW150914 allows us
		for the first time to constrain the coe cients in the PN series
		of the phasing up to 3.5PN order.
		\par 单参数分析中，一些参数的GR值落在后验分布对应的$2 \sim 2.5\sigma$，not surprising to the majority for corrrelation；
		进一步，如果有个noise导致某个参数的后验分布偏离0、其它所有参数也会同步偏离 (偏离方向对应每个PN系数符号)。
		\par 多参数分析中，PN系数由于相互的correlations、GR值都接近，除了0和0.5PN、对大多数 early-inspiral 的系数不敏感；
		而后面两个阶段，由于GW150914大部分SNR集中在这里，the constraints on those coe cients are very
		robust and essentially independent of the analysis configuration
		chosen, single or multiple.
		\par 最后，各个系数的log Bayesian factor分析表明no evidence in favor of the hypothesis that GR is violated
		(但是光这样的贝叶斯银子分析似乎不够充分、考虑到先验定义中人为设置的偏离宽度，Occam factor，见论文中的脚注7)。
	\item[2. ] GW170817\cite{test_GR_GW170817}
		\par 除了3.03.5PN，偏离为0都落在90\%可信区间里；
		\par 注意GW170817和GW150914的大多数PN系数偏离0点的方向都正好想反，后几阶可信区间宽度数量级差不多、前几阶好很多；

	\item[] 
\end{itemize}

\subsubsection{Dispersion Relation Test}


\subsubsection{Polarization Test}


\subsection{测量特定理论参数}
\par \cite{test_GR_GW170817}中提到对一些理论的约束，但是这些并没有参数化post-inspiral


\begin{itemize}
	\item[1. ] GW150914\cite{test_GR_GW150914}: 
	\item[2. ] GW170817\cite{test_GR_GW170817}:
		\par 对dipole parameter $B$ 的一些约束论述，详见论文；
	\item[] 
\end{itemize}


\newpage
\section{贝叶斯估计结果分析}
\subsection{Model Comparison}
\subsubsection{Bayes Factor and Odds Ratio}

\subsubsection{TIGER}

\subsection{Visualisation}


\subsubsection{Violin Plot}

\subsubsection{Re-express}


\newpage
\section{nonlinear conservative and dissipative effects}
\par \cite{test_GR_GW150914}：Those
 e ects include tails of radiation due to backscattering of GWs
 by the curved background around the coalescing black holes
 [75], nonlinear tails (i.e., tails of tails) [76], couplings between
 black-hole spins and the binary’s orbital angular momentum,
 interactions between the spins of the two bodies [77–79], and
 excitations of QNMs [28–30] as the remnant black hole settles
 in the stationary configuration.





\newpage
\bibliographystyle{plain}
\bibliography{ref}


\end{document}